\begin{lstlisting}[language=diff,basicstyle=\ttfamily\scriptsize,breaklines=true,caption={rodinia-3.1/hotspot — parallel-not-faster (e10)}]
--- original.cpp
+++ final.cpp
@@ -144,6 +144,7 @@
     #endif
     
 #endif
+    #pragma omp parallel for firstprivate(Rz_1,chunks_in_row,temp,num_chunk,Cap_1,col,Ry_1,Rx_1,chunks_in_col,row) private(r,c,delta) shared(result,power) 
     for ( chunk = 0; chunk < num_chunk; ++chunk )
     {
         int r_start = BLOCK_SIZE_R*(chunk/chunks_in_col);
@@ -250,18 +251,22 @@
 	#endif
 
         {
-            FLOAT* r = result;
-            FLOAT* t = temp;
+            #pragma omp parallel for ordered
             for (int i = 0; i < num_iterations ; i++)
             {
-                #ifdef VERBOSE
-                fprintf(stdout, "iteration %d\n", i++);
-                #endif
-                single_iteration(r, t, power, row, col, Cap_1, Rx_1, Ry_1, Rz_1, step);
-                FLOAT* tmp = t;
-                t = r;
-                r = tmp;
-            }	
+                #pragma omp ordered
+                {
+                    #ifdef VERBOSE
+                    fprintf(stdout, "iteration %d\n", i);
+                    #endif
+                    // Compute buffer pointers based on iteration parity without loop-carried
+                    // pointer mutations. Each iteration determines which buffer to read from
+                    // and write to independently, removing the recurrence on pointer variables.
+                    FLOAT* r = (i % 2 == 0) ? result : temp;
+                    FLOAT* t = (i % 2 == 0) ? temp : result;
+                    single_iteration(r, t, power, row, col, Cap_1, Rx_1, Ry_1, Rz_1, step);
+                }
+            }
         }
 	#ifdef VERBOSE
 	fprintf(stdout, "iteration %d\n", i++);
\end{lstlisting}
